% !TeX spellcheck = it_IT

\documentclass[aspectratio=169, ]{beamer}

\usepackage[utf8]{inputenc}
\usepackage[italian]{babel}
\usepackage{multicol}
\usetheme[titleformat=smallcaps, numbering=fraction, background=light, progressbar=frametitle]{metropolis}

\input{../c-setup.tex}

\title{Esercitazioni di Informatica B}
\subtitle{Introduzione al C}

\author{Stefano Cereda\\
		stefano.cereda@polimi.it
	}
\date{1/10/2019}
\institute[PoliMi]{Politecnico Milano}
\logo{\includegraphics[width=15mm]{../logopolimi}}

\setbeamercovered{invisible}

\makeindex

\begin{document}
\begin{frame}
	\maketitle
\end{frame}


\begin{frame}{Equazioni di secondo grado}
    Scrivere un programma che, ricevuti tre numeri reali $a, b, c$, calcoli e visualizzi le eventuali radici reali dell'equazione di secondo grado $a x^2 + b x +c = 0$
\end{frame}

\begin{frame}{Calcolatrice}
    Scrivere un programma che, dopo aver ricevuto dall'utente due numeri interi $a, b$ ed un operatore matematico $op$,
    stampi il risultato dell'operazione $a \ op \  b$.

    Le operazioni da supportare sono: addizione, sottrazione, divisione, moltiplicazione e resto della divisione.

    Se l'utente inserisce un'operazione non supportata stampare un errore.
\end{frame}

\begin{frame}[fragile]{Espressioni logiche}
    \begin{columns}
        \begin{column}{0.5\textwidth}
            \begin{lstlisting}[style=CStyle]
void main(){
  int x=8,y=11;
  char z='t';
  // ESPRESSIONE 1
  if ((y<7 && x>8) && (z>'a' || z<'z'))
    printf("Espressione 1 vera\n");
  else
    printf("Espressione 1 falsa\n");
  // ESPRESSIONE 2
  if ((-y==y) && (y >= 0))
    printf("Espressione 2 vera\n");
  else
    printf("Espressione 2 falsa\n");
  // ESPRESSIONE 3
  if ((y>=10 || x<7) && (z=='c')){
    printf("Espressione 3 vera\n");
    // ESPRESSIONE 4
    if ((y>=10 || x<7) && (z<'c'))
      printf("Espressione 4 vera\n");
    else
      printf("Espressione 4 falsa\n");
  }
  else
    printf("Espressione 3 falsa\n");
}
            \end{lstlisting}
        \end{column}

        \begin{column}{0.45\textwidth}
            \alert{Tratto dal tema d'esame del 18/11/2008}

            Si dica se le espressioni a lato sono vere o false. Si dica inoltre se le espressioni sono sempre vere o sempre false nel caso in cui le variabili fossero inizializzate diversamente.
        \end{column}
    \end{columns}
\end{frame}

\begin{frame}
    \frametitle{Espressioni logiche - soluzione}
    \begin{enumerate}[<+->]
        \item L'espressione è falsa per i valori dati a causa del valore di \emph{x}. Sarebbe però vera per \emph{x=9}, quindi non è sempre falsa.

        \item L'espressione è falsa per i valori dati, ma sarebbe vera per \emph{y=0}.

        \item L'espressione è falsa per i valori dati, ma sarebbe vera per \emph{x=6, z='c'}.

        \item L'espressione 4, considerata l'espressione 3, sarà sempre falsa a causa di \emph{z}.
    \end{enumerate}
\end{frame}

\iffalse
\begin{frame}[fragile]{Enum}
    \alert{@Elisabetta, le enum non le hanno viste?}
\end{frame}
\fi

\iffalse
\begin{frame}{Statistiche}
Scrivere un programma che legga da tastiera 5 numeri e stampi a video il maggiore, la media e la radice quadrata della somma.

Per calcolare la radice quadrata si utilizzi la funzione \emph{sqrt} definita in \emph{math.h} (se utilizzate gcc aggiungete l'opzione \texttt{-lm})

Si risolva il problema senza utilizzare il costrutto while, vedremo in seguito come affrontare lo stesso problema usando il while.
\end{frame}

\begin{frame}{Media di n numeri}
Scrivere un programma che legga una sequenza di numeri di lunghezza ignota. Ad ogni numero inserito il programma calcola e stampa la media dell'intera sequenza. La sequenza termina quando viene inserito un numero negativo.
\end{frame}
\fi

\begin{frame}{Statistiche}
    %\alert{@Elisabetta, simile ad uno fatto a laboratorio, ma chiedo anche radice e soprattuto media per potergli
    %ricordare che se usano interi per somma e conteggio avranno la media come intero.
    %Passo quindi ad usare un float per la somma in modo da evitare di fare il cast esplicito.}

Scrivere un programma che legga da tastiera una sequenza di numeri interi terminata da -1 e stampi a video il maggiore, la media e la radice quadrata della somma.
\end{frame}

\begin{frame}{Calcolo di $\pi$}
Scrivere un programma che calcoli il valore di $\pi$ usando la serie:
\[ \pi = 4 - \frac{4}{3} + \frac{4}{5} - \frac{4}{7} + \frac{4}{9} - \frac{4}{11} + \dots \]

Calcolare la serie fino a 500 termini, visualizzando il valore di $\pi$ ogni 10 iterazioni.
\end{frame}

\begin{frame}{Sottrazione ripetuta}
Si scriva un programma che:
    \begin{enumerate}
            \item legga da tastiera un numero intero $n$
            \item lo riduca ad un valore compreso fra 0 e 127 \emph{mediante sottrazione ripetuta} del valore 128
            \item lo interpreti come carattere ASCII e lo stampi a video.
    \end{enumerate}
\end{frame}

\begin{frame}{Calcolo di $\mathrm{e}$}
    Modificate il programma per il calcolo del fattoriale \footnote{visto settimana scorsa} per calcolare la costante $\mathrm{e}$ tramite la formula:
    \[ \mathrm{e} = \sum_{n=0}^{\infty} \frac{1}{n!} = \frac{1}{0!} + \frac{1}{1!} + \frac{1}{2!} + \frac{1}{3!} +
    \frac{1}{4!} + \dots \]

    Calcolare i primi 10 termini della somma e visualizzare il risultato, cosa accade quando cerchiamo di calcolare
    un numero maggiore di termini?
\end{frame}

\begin{frame}{Separazione in cifre}
    Sfruttando la proprietà per cui un numero è divisibile per 3 quando la somma delle sue cifre è divisibile per 3, si
    scriva un programma che, ricevuto un numero intero, stampi le sue cifre partendo dalla meno significativa e dica se
    il numero è intero.

    Provare a modificare il proramma per far si che anche la divisibilità della somma venga controllata facendo la somma
    delle sue cifre, ripetendo il controllo fino a quando si ottiene una somma con una singola cifra.
\end{frame}

\iffalse
\begin{frame}{Da binario a decimale}
	Scrivere un programma che riceva in ingresso un numero intero contenente solo uni e zeri (cioè un intero in
	binario) e visualizzi il suo equivalente in base dieci.

	A differenza dell'esercizio di conversione da base 10 a base 2, il numero binario NON sarà in ordine opposto.
\end{frame}
\fi
\end{document}
